\section{Discussion}\label{sec:discussion}

In the DR-grading settings studied here, an external failure separated into two components that
call for different remedies. The ranking of images by severity largely transferred to new
populations, also for models whose final training set was EyePACS rather than the development
pool; the mapping from score to grade did not, and its error had opposite signs on different
populations because the score given the true grade moved. Where ranking survived, part of the loss was a property of the decision rule and
could be recovered without retraining --- the mild grade on APTOS, referral sensitivity in the
reverse direction --- but only as a trade between grades, and only with local labels. Where ranking
degraded, as for the foundation model on DDR, no rule helped. Two of our own predictions bound the
argument: deployable rules did not outweigh the choice of model, and the collapse of mild recall
that the literature reads as representational blindness was specific to one corpus, although mild
F1 stayed low on every population. A reader who sees
an external failure should therefore ask, in order, whether the model still ranks the new
population, where its cut-points were set and on which population, and only then what it has
learned.

These results refine rather than contradict recent work. ORDER-DR separated threshold-dependent
grade assignment from threshold-independent risk ranking and noted threshold shift on external
data\cite{sheng2026orderdr}; reanalysing its released predictions shows that its reported
advantage over its own branches does not survive the matched decision rule. Tabibu and Abhyankar
found, in a preprint, that the probability calibration of a binary referral model degrades
across cameras mostly through prevalence, and that 100--200 labelled images restore
it\cite{tabibu2026discrimination}. The mechanism we found differs: the score given the true grade
moved, so prior adaptation without labels failed even against its own decoding, and
temperature-style rescaling cannot be expected to help either, as reported for frozen
encoders\cite{poyrazer2026calibration}. For a referral threshold, the number of labels needed is
set by the referable cases among them, and the risk follows a distribution-free form, so a few
hundred images suffice only where referable disease is common. A public, not yet peer-reviewed
model repository reports that per-threshold calibration degrades unevenly under domain shift and
that the mild grade fails on every held-out domain\cite{gradeye2026}, consistent with the
grade-specific movement and the low mild F1 we observed.

The central objection to this kind of analysis is that refitting cut-points on local labels may
fit local label noise rather than fix the model. The populations answer it differently. On DDR every
grade moved down by a similar amount, yet the images were not degraded on any measure we could
compute, although degradation can produce exactly this pattern; the change lies in what the images
show or in how they were graded, and recalibration corrects the mapping without revealing which. On
APTOS one grade moved by about a whole step while its neighbours did not, which a change in image
quality does not produce and a difference in how a single grader assigned the mild grade would;
recalibration there aligns the model with the local grading standard. The models trained on EyePACS
add a second line of evidence: they under-graded DDR as the development-pool models did, yet their
EyePACS-fitted threshold over-referred there, because a transferred threshold carries the reference
standard of the set it was fitted on. Where a threshold lands is thus partly a property of the
labels on both sides, which is consistent with, though not proof of, a difference in grading
convention; re-grading a random sample of APTOS mild and DDR images by one ophthalmologist would
test it directly. Whether recalibration is a fix depends on whether the local standard is the one
the screening programme acts on --- which, for a programme deploying a model, it usually is. Either
way, the rule must be fitted where it will be used, on a sample large enough for its risk to be
known.

The practical recommendations cost no additional training. External evaluations should report
discrimination at each ordinal threshold beside grade-level agreement, the full referral operating
curve, and the transfer gap of any development-fitted threshold, since the same threshold landed
\pn{aptos.tr.e08.A.gap} and \pn{ddr.tr.e08.A.gap} points from an 80\% target on two populations.
Model comparisons should be made under a matched decision rule that names its fitting procedure,
with differences judged against the variation between seeds of one configuration. Deployment
should include setting the referral threshold on a local sample, sized by its referable cases
rather than its images: about 120 referable cases keep the probability of missing an 80\%
sensitivity floor near one in ten and about 200 near one in twenty, which at a referable prevalence
of 10\% means 1\,000 to 2\,000 consecutive images; where fewer cases are available, setting the cut
for a higher sample sensitivity buys the same safety with specificity. This is in line with advice
on updating prediction models for new
settings\cite{janssen2008updating,vergouwe2017updating,riley2021samplesize}, and because the risk
does not depend on the model or the population it can be planned before any local data are
collected.

The study has limitations. All data are public and retrospective, the populations are defined by
dataset rather than by a prospective screening programme, and APTOS and EyePACS labels come from
single graders. SUSTech-SYSU is a case--control collection with few mild images, and APTOS's no-DR
and mild images separate unusually easily, partly through acquisition, so discrimination there
overstates what a screening population would show; acquisition statistics also carried grade
information within the development pool, whose two corpora differ in both storage and grade mix
(estimated JPEG quality, area under the curve \pn{acq.dev.jpegquality.ref}), a shortcut the models
may have used. Patient identifiers are unavailable for most corpora, so resampling is by image
and fellow-eye pairs may be correlated. The external-label rules (R4, R4m) are diagnostics of what
local labels could recover, not external validation; the deployable recommendation rests on the
separate sample-size analysis. Referable disease was defined by DR grade alone because macular
oedema is not graded in the external sets, whereas real screening refers for both; the
macular-oedema head could not be validated externally and is reported only in the Supplement. The
reverse direction, the seeds and the large-data models each use one or two architectures, and the
EyePACS models selected their checkpoint on the same fold that set their development threshold.
Finally, the analysis plans were committed to version control rather than to a public registry,
the pre-specification of development runs trained before the repository history begins cannot be
verified, the second-stage prediction on the direction of the shift was written knowing it for
APTOS and DDR, and the analyses added after an internal editorial assessment, including the
sample-size analysis in referable cases, are post hoc. The decision-rule effects should be confirmed prospectively, with cut-points fitted on
an initial local sample and then frozen.

\section{Conclusions}\label{sec:conclusions}

External failures and apparent differences between DR grading models are routinely attributed to
learned representations. In the settings studied here a substantial part of them lay in the
decision rule that converts scores into grades and could be diagnosed and repaired without
retraining, while the remainder reflected real losses of ranking that no rule recovers. Reporting
per-threshold discrimination, full operating curves, the transfer gap of thresholds and
matched-decision-rule comparisons, and setting referral thresholds on a local sample with enough
referable cases before deployment, would separate the two.
